\section{Ordered Requests: Mention Order Coordinates Slots, Not Agents}
\label{sec:ordered}

\begin{table}[t]
\caption{Committing more tokens per step, with exact slot lengths (BFCL parallel and parallel\_multiple, 400 requests). Set accuracy and \ccer in \%; \ccer{} (sym.) on the requests whose slot lengths cannot tell the calls apart (213 for \dream, 216 for \llada); order: share of symmetric \texttt{parallel} requests whose $i$-th call took the $i$-th mentioned entity; fwd: forward passes per request. LTR: left to right, one token per step; $\tau$: commit every position above confidence 0.9. \qwen (autoregressive) in the same skeleton: 86.5\% set accuracy, 1.3\% \ccer.}
\label{tab:ksweep}
\small\setlength{\tabcolsep}{3.2pt}
\begin{tabular}{@{}lrrrrrrr@{}}
\toprule
 & $k{=}1$ & 2 & 4 & 8 & 16 & LTR & $\tau$ \\
\midrule
\multicolumn{8}{@{}l}{\emph{\dream (full attention)}} \\
Set acc. & 89.8 & 88.0 & 87.0 & 86.3 & 84.5 & 89.3 & 89.0 \\
\ccer & 1.3 & 1.3 & 1.5 & 1.8 & 2.0 & 1.0 & 1.3 \\
\ccer (sym.) & 1.9 & 1.9 & 2.3 & 2.8 & 2.8 & 1.4 & 1.9 \\
Order & 95.2 & 95.2 & 95.0 & 96.1 & 96.1 & 97.5 & 95.1 \\
Fwd & 28.0 & 14.2 & 7.3 & 3.9 & 2.2 & 27.8 & 1.9 \\
\midrule
\multicolumn{8}{@{}l}{\emph{\llada (32-token blocks)}} \\
Set acc. & 87.8 & 87.5 & 86.5 & 83.3 & 79.8 & -- & -- \\
\ccer & 2.5 & 2.8 & 3.3 & 4.0 & 4.5 & -- & -- \\
\ccer (sym.) & 0.9 & 0.9 & 1.4 & 1.4 & 1.9 & -- & -- \\
Order & 95.1 & 96.3 & 96.3 & 95.1 & 93.8 & -- & -- \\
Fwd & 27.8 & 14.8 & 8.3 & 5.1 & 3.9 & -- & -- \\
\bottomrule
\end{tabular}
\end{table}

\begin{table}[t]
\caption{One canvas against teams of agents. Ordered: set accuracy (\%) on the 104 BFCL requests whose sibling calls share one single-call skeleton. Symmetric (choose-N): Dup.\ is the share of requests (\%) in which two calls take the same city; Order, the share of list requests whose $i$-th call takes the $i$-th listed city. Agents act at the same time or in turns (seeing the calls made so far) and are anonymous, numbered (``assistant $i$ of $n$''), or numbered and told that assistant $i$ makes the $i$-th call. Greedy decoding on 60 list and 45 open requests; $T{=}0.7$: five seeds. Open requests hold one-token slots, so only duplicates are reported.}
\label{tab:teams}
\small\setlength{\tabcolsep}{2.6pt}
\begin{tabular}{@{}lrrrrrr@{}}
\toprule
 & Ordered & \multicolumn{3}{c}{Symmetric, greedy} & \multicolumn{2}{c}{$T{=}0.7$} \\
\cmidrule(lr){2-2}\cmidrule(lr){3-5}\cmidrule(l){6-7}
 & Set & \multicolumn{2}{c}{List} & Open & List & Open \\
\cmidrule(lr){3-4}
 & acc. & Dup. & Order & Dup. & Dup. & Dup. \\
\midrule
\multicolumn{7}{@{}l}{\emph{One canvas holds all $n$ calls}} \\
\dream, $k{=}16$ & 81.7 & 15.0 & 11.7 & 100.0 & 17.3 & 100.0 \\
\dream, $\tau{=}0.9$ & 86.5 & 0.0 & 48.3 & 0.0 & -- & -- \\
\dream, $k{=}1$ & 88.5 & 0.0 & 48.3 & 0.0 & 0.0 & 0.4 \\
\llada, $k{=}16$ & 81.7 & 0.0 & 96.7 & 4.4 & -- & -- \\
\llada, $k{=}1$ & 88.5 & 0.0 & 100.0 & 0.0 & -- & -- \\
\qwen (AR) & 91.3 & 0.0 & 81.7 & 0.0 & -- & -- \\
\midrule
\multicolumn{7}{@{}l}{\emph{A team of $n$ \qwen agents, one call each}} \\
\quad same time, anonymous & 0.0 & 100.0 & 0.0 & 100.0 & 99.7 & 93.8 \\
\quad same time, numbered & 0.0 & 100.0 & 0.0 & 97.8 & 100.0 & 93.8 \\
\quad same time, told the rule & 34.6 & 70.0 & 8.3 & -- & -- & -- \\
\quad turns, anonymous & 88.5 & 0.0 & 91.7 & 0.0 & -- & -- \\
\quad turns, numbered & 88.5 & 0.0 & 91.7 & 0.0 & -- & -- \\
\midrule
\multicolumn{7}{@{}l}{\emph{A team of $n$ \dream agents, one call each}} \\
\quad same time, anonymous & 0.0 & 100.0 & 0.0 & 100.0 & 87.7 & 100.0 \\
\quad same time, numbered & 28.8 & 78.3 & 3.3 & 100.0 & 75.0 & 100.0 \\
\quad same time, told the rule & 48.1 & 45.0 & 10.0 & -- & -- & -- \\
\quad turns, anonymous & 84.6 & 0.0 & 35.0 & 0.0 & -- & -- \\
\quad turns, numbered & 85.6 & 0.0 & 26.7 & 0.0 & -- & -- \\
\bottomrule
\end{tabular}
\end{table}


We first test \pred{1} on the BFCL requests with exact slot lengths (Table~\ref{tab:ksweep}).

\paragraph{Same-step slots coordinate.} Committing 16 tokens per step instead of one lowers \dream's set accuracy from 89.8\% to 84.5\% (by 5.3 points, 95\% interval [3.3, 7.5]) and \llada's from 87.8\% to 79.8\%, while cutting the forward passes per request by factors of 13 and 7. The confidence threshold keeps \dream at 89.0\% with 1.9 forward passes per request; the 0.8-point loss is not significant ($p=0.45$). Cross-call errors stay rare: \dream's rate does not change significantly between $k=1$ and $k=16$ (1.3\% and 2.0\%, $p=0.38$), \llada's rises from 2.5\% to 4.5\% ($p=0.02$), and the AR agent's is 1.3\%. On the equal-length requests, where slot lengths cannot tell the calls apart, \dream's set accuracy falls from 86.9\% to 81.7\%, much as on all requests. On those in the \texttt{parallel} category, the $i$-th call takes the $i$-th mentioned entity in 94.0\% of \dream's requests with one token per step and in 88.1\% with 16; among the outputs that parse with every call matched, the share stays between 95.0\% and 96.1\% at every $k$. The cost of parallelism lies inside single values: of the 21 requests that \dream gets right with one token per step and wrong with 16, 13 have only single-call errors and 5 no longer parse, against 3 with cross-call errors (\llada: 19, 7, and 7 of 33). BFCL usually lists the reference calls in mention order and our skeleton lays them out in that order, so following the order is also correct; and most values are copied from the request with high confidence, so the slots rarely face a real choice.

\paragraph{Simultaneous agents do not.} On the 104 team-symmetric requests, the agents that call the same tool see the same skeleton, so only their prompts can set them apart (Table~\ref{tab:teams}, left column). Anonymous agents that act at the same time receive identical inputs and make identical calls, so no team is correct. Numbering them helps \dream's agents a little and \qwen's not at all: 28.8\% and 0\% set accuracy, against 81.7\% for \dream's canvas at 16 tokens per step on the same requests (difference 52.9 points, [42.3, 62.5]). Even when told the convention, they reach only 48.1\% and 34.6\%. Most numbered agents take the entity mentioned first (67.8\% of \dream's and 88.3\% of \qwen's agents in sibling groups). Agents that take turns come close to the single agent but stay below it: \dream's teams reach 84.6--85.6\% against 88.5\% for its canvas with one token per step, and \qwen's 88.5\% against 91.3\% for its single agent; each team loses 3 or 4 requests that the single agent gets right and gains at most one ($p\geq0.125$).

\pred{1} holds for same-step slots: a position on the shared skeleton acts as a role, and the request's order assigns the entities to roles. It fails for agents that act at the same time, for whom a number in the prompt, and even a stated rule, is a weak role. Taking turns comes within 4 points of the single agent.
